
\medskip
\phantomsection\label{par:sga3-VIA-4.3}%
\noindent\textbf{4.3.---}
It follows from the preceding discussion that the groupoid \(G_*\),
with base \(G\), defined by \(u\), is the direct sum of the groupoids
\(\mathscr S(G^\alpha_*)\) induced by \(G_*\) on the various
open-and-closed parts of \(G\) of the form
\(\mathscr S(G^\alpha)\). The cokernel of \(G_*\) is therefore the
direct sum of the cokernels of the groupoids
\(\mathscr S(G^\alpha_*)\), which we are thus led to study
separately.

First consider the groupoid \(\mathscr S(G^0_*)\) induced by \(G_*\)
on \(\mathscr S(G^0)\). It is clear that
\(\mathscr S(G^0_*)\) is the groupoid with base
\(\mathscr S(G^0)\) defined by the homomorphism
\[
  F\longrightarrow\mathscr S(G^0)
\]
induced by \(u\) (\S~\hyperref[par:sga3-VIA-3.1]{3.1}). The cokernel whose existence we wish to
prove is thus identified with \(F\backslash\mathscr S(G^0)\).
Consider, on the other hand, the groupoid
\begin{center}
  \phantomsection\label{diag:sga3-VIA-4.3-G0-groupoid}%
  \begin{tikzcd}[column sep=large]
    G^0_2
      \arrow[r,shift left=1.4ex,"\ell'_2"]
      \arrow[r,"\ell'_1"]
      \arrow[r,shift right=1.4ex,"\ell'_0"'] &
    G^0_1
      \arrow[r,shift left=0.8ex,"\ell_1"]
      \arrow[r,shift right=0.8ex,"\ell_0"'] &
    G^0_0=G^0
  \end{tikzcd}
\end{center}
induced by \(\mathscr S(G^0_*)\) on \(G^0\). Referring to the
explicit construction in Expos\'e~V, \S~3.b, the object denoted there
by \(Y_0\times_{X_0}X_1\) is precisely \(F\times_A G^0\). Hence
\(G^0_1\) is the inverse image of \(G^0\) under the morphism
\[
  F\times_A G^0\longrightarrow\mathscr S(G^0)
\]
induced by \(\lambda\).

I claim that this inverse image is \(F_0\times_A G^0\), where
\(F_0\) denotes the inverse image of \(G^0\) under \(u\). Indeed, if
\(F^\beta\) is a connected component of \(F_0\), then
\(F^\beta\times_A G^0\) is connected by \SGARefLink{lem:sga3-VIA-2.1.2}{2.1.2} and
\(\lambda(F^\beta\times_A G^0)\) is contained in \(G^0\).%
\begingroup
\renewcommand{\thefootnote}{\fnsymbol{footnote}}%
\footnote[1]{Translator's note: the French source prints
\(\lambda(F\times_A G^0)\) at this point. The statement concerns the
connected component \(F^\beta\) just introduced, and the converse
immediately below confirms the intended restriction to
\(F^\beta\times_A G^0\).}%
\endgroup
Conversely, if \(F^\beta\) is a connected component of \(F\) not
contained in \(F_0\), then the image of
\(F^\beta\times_A G^0\) is again connected and contains
\(u(F^\beta)\); if \(u(F^\beta)\) is not contained in \(G^0\), then
\(\lambda(F^\beta\times_A G^0)\) does not meet \(G^0\).

It follows that the groupoid \(G^0_*\) induced by \(G_*\) on \(G^0\)
is the groupoid with base \(G^0\) defined by the homomorphism
\[
  F_0\longrightarrow G^0
\]
induced by \(u\). Since \(G^0\), and hence \(F_0\), are of finite
type over \(A\), \hyperref[sec:sga3-VIA-3]{Section~3} shows that \(G^0_*\) has a cokernel,
namely \(F_0\backslash G^0\).%
\begingroup
\renewcommand{\thefootnote}{\fnsymbol{footnote}}%
\footnote[2]{Translator's note: the French source says
``by \hyperref[sec:sga3-VIA-4]{paragraph 4}.'' The finite-type quotient invoked here is the
result of \hyperref[sec:sga3-VIA-3]{Section~3}; \hyperref[sec:sga3-VIA-4]{Section~4} is the general-case construction
currently in progress.}%
\endgroup

I now claim that \(F_0\backslash G^0\) is identified with
\(F\backslash\mathscr S(G^0)\). The proof is analogous to that of
the first part of assertion~(i) of \SGARefLink{lem:sga3-V-6.1}{Expos\'e~V, Lemma~6.1}. Consider
the diagram
\begin{center}
  \phantomsection\label{diag:sga3-VIA-4.3-span}%
  \begin{tikzcd}[column sep=large]
    \mathscr S(G^0) &
    F\times_A G^0
      \arrow[l,"v"']
      \arrow[r,"\operatorname{pr}_2"] &
    G^0
  \end{tikzcd},
\end{center}
where \(v\) is the morphism induced by \(\lambda\). Since
\(\operatorname{pr}_2\) has a section, it is a universal effective
epimorphism. Consequently, \(F_0\backslash G^0\) coincides with
\(\operatorname{Coker}(v_0,v_1)\), where
\begin{center}
  \phantomsection\label{diag:sga3-VIA-4.3-V-groupoid}%
  \begin{tikzcd}[column sep=large]
    V_2
      \arrow[r,shift left=1.4ex,"v'_2"]
      \arrow[r,"v'_1"]
      \arrow[r,shift right=1.4ex,"v'_0"'] &
    V_1
      \arrow[r,shift left=0.8ex,"v_1"]
      \arrow[r,shift right=0.8ex,"v_0"'] &
    V=F\times_A G^0
  \end{tikzcd}
\end{center}
is the inverse image under \(\operatorname{pr}_2\) of the groupoid
\(G^0_*\) (cf.\ Expos\'e~V, \S~3.a), and hence also the inverse image
of \(\mathscr S(G^0_*)\) under the composite morphism
\begin{center}
  \phantomsection\label{diag:sga3-VIA-4.3-composite}%
  \begin{tikzcd}[column sep=large]
    F\times_A G^0
      \arrow[r,"\mathrm{inclusion}"] &
    F\times_A\mathscr S(G^0)
      \arrow[r,"\operatorname{pr}_2"] &
    \mathscr S(G^0)
  \end{tikzcd}.
\end{center}

Similarly, since \(v\) is faithfully flat and quasi-compact,
\(F\backslash\mathscr S(G^0)\) coincides with the cokernel of the
inverse image of \(\mathscr S(G^0_*)\) under the base change \(v\).
This inverse image is isomorphic to \(V_*\) by Expos\'e~V, \S~3.c.
It follows that the canonical inclusion of \(G^0_*\) in
\(\mathscr S(G^0_*)\) induces an isomorphism
\[
  F_0\backslash G^0
    \xrightarrow{\ \sim\ }
  F\backslash\mathscr S(G^0).
\]

Finally, formation of \(F\backslash\mathscr S(G^0)\) commutes with
finite locally free base changes, because the same holds for
\(F_0\backslash G^0\) (cf.\ editor's
note~\ref{note:sga3-VIA-37} and \SGARefLink{par:sga3-VIA-4.6}{4.6} below).
